\chapter*{Suggested Reading}
\addcontentsline{toc}{chapter}{Suggested Reading}
\markboth{Suggested Reading}{Suggested Reading}

This atlas bridges the geometry of inscribed convex polygons with permutation group actions and circular composition combinatorics. The curated works below provide foundational theory, classical historical developments, and broader algebraic context.

\section*{Dihedral Enumeration and Group Actions}

\begin{description}
    \item[\textbf{Augustin-Louis Cauchy (1845)}]\mbox{}\\
    \emph{M\'emoire sur diverses propri\'et\'es remarquables des substitutions r\'eguli\`eres ou irr\'eguli\`eres, et des syst\`emes de substitutions conjugu\'ees}, \emph{Comptes Rendus de l'Acad\'emie des Sciences}, 21:972--987.\\
    The original formulation of the orbit-counting lemma for permutation groups.

    \item[\textbf{Ferdinand Georg Frobenius (1887)}]\mbox{}\\
    \emph{\"Uber die Congruenz nach einem aus zwei endlichen Gruppen gebildeten Doppelmodul}, \emph{Journal f\"ur die reine und angewandte Mathematik (Crelle's Journal)}, 101:273--299.\\
    The generalisation of the orbit-counting theorem to abstract finite group actions.

    \item[\textbf{William Burnside (1911)}]\mbox{}\\
    \emph{Theory of Groups of Finite Order}, 2nd edition, Cambridge University Press, Cambridge.\\
    The classical treatise popularising the orbit-counting lemma (attributing the theorem to Frobenius in Section~118), forming the arithmetic core of dihedral cycle index evaluations.

    \item[\textbf{Peter M. Neumann (1979)}]\mbox{}\\
    \emph{A story that is not Burnside's}, \emph{The Mathematical Gazette}, 63(426):271--270.\\
    A delightful historical account detailing the provenance of the Cauchy--Frobenius lemma and its popular misattribution to Burnside.

    \item[\textbf{J. Howard Redfield (1927)}]\mbox{}\\
    \emph{The theory of group-reduced distributions}, \emph{American Journal of Mathematics}, 49(3):433--455.\\
    The pioneering paper anticipating P\'olya's enumeration theorems, introducing group-reduced distributions and the superposition of permutation representations.

    \item[\textbf{George P\'olya (1937)}]\mbox{}\\
    \emph{Kombinatorische Anzahlbestimmungen f\"ur Gruppen, Graphen und chemische Verbindungen}, \emph{Acta Mathematica}, 68(1):145--254.\\
    The foundational memoir formalising the cycle index polynomial $Z(G; z_1, \dots, z_n)$ and its functional substitution principles for weighted configuration counting under group equivalence.

    \item[\textbf{Adalbert Kerber (1999)}]\mbox{}\\
    \emph{Applied Finite Group Actions}, 2nd edition, Springer-Verlag, Berlin, Heidelberg.\\
    An exhaustive modern reference on the algebraic and algorithmic apparatus of permutation representations, double-coset decompositions, and constructive orbit generation.
\end{description}

\section*{Cyclic Polygons and Circular Compositions}

\begin{description}
    \item[\textbf{Duncan M. Y. Sommerville (1909)}]\mbox{}\\
    \emph{On Certain Periodic Properties of Cyclic Compositions of Numbers}, \emph{Proceedings of the London Mathematical Society}, Second Series, 7:263--313.\\
    A fundamental early 20th-century investigation into the period structures and symmetry classes of cyclic integer compositions, establishing the combinatorial basis for circular stride sequences.

    \item[\textbf{Hansraj Gupta (1979)}]\mbox{}\\
    \emph{Enumeration of Incongruent Cyclic $k$-gons}, \emph{Indian Journal of Pure and Applied Mathematics}, 10(8):964--999.\\
    An exact combinatorial classification of incongruent convex polygons inscribed in circular lattices under physical rotation and reflection.

    \item[\textbf{E. N. Gilbert and John Riordan (1961)}]\mbox{}\\
    \emph{Symmetry Types of Periodic Sequences}, \emph{Illinois Journal of Mathematics}, 5(4):657--665.\\
    The definitive classification of periodic circular sequences (necklaces and bracelets) under cyclic and dihedral group actions.

    \item[\textbf{Petros Hadjicostas (2017)}]\mbox{}\\
    \emph{Cyclic, Dihedral and Symmetrical Carlitz Compositions of a Positive Integer}, \emph{Journal of Integer Sequences}, 20(8):Article 17.8.5.\\
    A modern treatment of cyclic and dihedral composition spaces with restricted adjacent summands and bilateral symmetry constraints.
\end{description}

\section*{Combinatorial and Graphical Enumeration}

\begin{description}
    \item[\textbf{Richard P. Stanley (2011)}]\mbox{}\\
    \emph{Enumerative Combinatorics}, Volume 1, 2nd edition, Cambridge University Press, Cambridge.\\
    The standard modern reference on generating functions, posets, and group actions on finite sets, including necklace and bracelet algebra.

    \item[\textbf{Frank Harary and Edgar M. Palmer (1973)}]\mbox{}\\
    \emph{Graphical Enumeration}, Academic Press, New York.\\
    A classic, pedagogical text applying P\'olya enumeration techniques to geometric graphs, simplicial complexes, and spatial symmetry classes.

    \item[\textbf{George E. Andrews (1976)}]\mbox{}\\
    \emph{The Theory of Partitions}, Addison-Wesley, Reading, MA.\\
    The authoritative treatise on integer partitions, modular identities, and their geometric realizations.
\end{description}

\section*{Geometric Metric Spaces and Musical Orbifolds}

\begin{description}
    \item[\textbf{Godfried T. Toussaint (2013)}]\mbox{}\\
    \emph{The Geometry of Musical Rhythm: What Makes a ``Good'' Rhythm Good?}, CRC Press, Boca Raton, FL.\\
    A comprehensive geometric investigation into cyclic rhythmic timelines, analysing maximal evenness, convex hulls, and Afro-Cuban clave patterns as inscribed polygons in $\mathbb{Z}_{16}$.

    \item[\textbf{Dmitri Tymoczko (2011)}]\mbox{}\\
    \emph{A Geometry of Music: Harmony and Counterpoint in the Extended Common Practice}, Oxford University Press, New York.\\
    Foundational formulation of musical scales, pitch-class sets, and voice leading as geometric configurations and paths within quotient orbifolds under dihedral and symmetric group actions.
\end{description}
